% FACT_PAYLOAD：以下内容均来自 supplied artifact，可用于起草。
2022 年公报口径把江豚 1249 头、长江鲟 438 尾作为问题二的初值与口径锚定。

\(\delta_B\) 表示通道阻隔和栖息地碎片化的综合损失；\(R_S\) 表示人工放流强度。

基线修复情景下，江豚和长江鲟的归一化末值分别为 1.133 和 1.140。

仅提高通道阻隔、不叠加污染时，江豚和长江鲟的末值分别为 0.901 和 0.856，相对基线分别下降 20.5\% 和 24.9\%。

将人工放流项置为 \(R_S=0\) 时，江豚末值为 1.145，长江鲟末值为 0.519；相对基线，江豚变化为 +1.0\%，长江鲟下降 54.5\%。这一对照用于拆分自然恢复与人工放流补偿效应。

污染情景中，江豚和长江鲟的末值分别为 0.623 和 0.706，相对基线分别下降 45.0\% 和 38.0\%。这一情景用于比较污染叠加阻隔对基线收益的侵蚀。

上述情景输出都是当前参数组与 2022 公报归一化口径下的内部投影，不构成独立外部验证。\(M\) 和 \(V\) 是传导和竞争状态，不是本问的直接观测目标。

% EDITORIAL_REQUIREMENT：写成“问题二结果与讨论”的两个自然段。第一段报告基线与三个对照的关键结果；第二段解释这些对照能够支持的机制判断与认识论边界。不要使用列表或小标题。
